% Section: Introduction
% =====================================================================
% DESCRIPTION
% -----------
% Goal: motivate the intensity-injection hypothesis, state the gap,
% declare a negative result honestly, and preview the Steam paradox so
% the reader knows the paper has a mechanism story (not just a null).
%
% Paragraph plan (~5 paragraphs):
%   P1. Sequential recommenders treat history as a token stream;
%       attention models (SASRec, BERT4Rec) dominate. But every
%       interaction is collapsed to a binary "happened" — the
%       underlying *intensity* (rating, hours played, dwell time) is
%       discarded. That is a lot of supervised signal thrown away.
%   P2. A natural hypothesis: re-introducing per-interaction intensity
%       into self-attention should help, in proportion to how rich the
%       intensity signal is. Steam playtime (continuous, heavy-tailed)
%       should help more than 5-star ratings (discrete, low entropy).
%   P3. We test this with three drop-in modifications to SASRec
%       (Add / Mul / Val), each adding ONE learnable scalar per layer
%       so that lambda=0 collapses exactly to SASRec — the model can
%       opt out if intensity is uninformative.
%   P4. Headline finding: the hypothesis fails. Across 7 datasets and
%       5 seeds per (dataset, model), no variant achieves a robust
%       NDCG@10 / Recall@10 improvement, and the failure is most
%       pronounced on Steam — the OPPOSITE of what the hypothesis
%       predicts. We confirm this with paired t-tests, per-user paired
%       bootstrap CIs, and per-seed sign-agreement.
%   P5. Mechanism: learned lambda stays high (0.55-0.99 on Steam) even
%       as metrics degrade, so the "opt-out" gate fails to fire. The
%       paper's contribution is the documented failure together with
%       the lambda-vs-Delta calibration evidence that simple
%       per-layer scalar gating is too coarse for heavy-tailed
%       intensities.
%
% Reusable prose: ia_sasrec.md §1 (motivation paragraphs) — lift but
% retune the tone from "we propose" to "we hypothesised".
%
% Contributions bullet list at end of section:
%   - Three IA-SASRec variants with explicit lambda fallback.
%   - A 7-dataset, 5-seed evaluation with a dual-criterion significance
%     protocol (paired t + per-user bootstrap + per-seed sign agreement).
%   - A negative result with mechanism analysis via the learned lambda.
%   - Released code, checkpoints, and per-user metric parquets for full
%     reproducibility.
% =====================================================================
\section{Introduction}
\label{sec:intro}

Sequential recommenders treat a user's interaction history as a token
stream, and self-attention models such as SASRec~\citep{kang2018sasrec}
and BERT4Rec~\citep{sun2019bert4rec} dominate the leaderboards. In
the standard formulation, every interaction is collapsed to a binary
``happened'': item $i_j$ was in the user's history at position $j$,
and that is the totality of the signal that reaches the attention
layer. The underlying intensity of each interaction (a
$5$-star rating on MovieLens, hours played on Steam, dwell time or
purchase value in industrial logs) is discarded. From a Shannon
standpoint that is a lot of supervised information thrown away, and
the discard is more notable the richer the intensity signal is:
a binary signal carries one bit; a continuous heavy-tailed playtime
distribution can carry several.

A natural hypothesis follows: re-introducing per-interaction intensity
into the attention computation should help, and should help in
proportion to how informative the underlying signal is. Steam
playtime (continuous, heavy-tailed, several bits per interaction)
should help more than discrete $1$--$5$ ratings on
MovieLens, whose entropy is bounded above by $\log_2 5 \approx 2.3$
bits and is much lower in practice given the well-known rating bias
toward $4$ and $5$.

We test this hypothesis with three drop-in modifications of SASRec
(IA-SASRec-Add, IA-SASRec-Mul, and IA-SASRec-Val), corresponding
to the three algebraic positions in the attention expression where
intensity can be threaded in (pre-softmax additive bias, pre-softmax
multiplicative scale, post-softmax attention reweighting).
Each variant adds exactly one learnable scalar $\lambda$ per
attention layer, initialised to $1.0$, designed so that $\lambda{=}0$
collapses the variant bit-for-bit to vanilla SASRec. The model is
therefore a strict super-set of the baseline: if intensity is
not informative, a fair learner can recover SASRec exactly.

Across seven datasets, spanning the discrete-rating and
continuous-playtime regimes, and five random seeds per
(dataset, model) cell, the hypothesis fails. Not one of the $63$
tested cells ($3$ variants $\times$ $3$ metrics at $k{=}10$ across
$7$ datasets) shows a robust improvement under our dual-criterion
significance protocol: a seed-level paired $t$ and a per-user
paired bootstrap that must both reject, guarded by a per-seed
CI sign-majority on the bootstrap. The single
positive paired-$t$ cell (IA-SASRec-Add on Amazon-Office MRR@10,
$+3.6\%, p{=}0.025$) is suppressed by the per-user bootstrap and
the per-seed sign agreement. Worse, the degradation pattern is the
opposite of the intensity-richness prediction. Steam, the
dataset with the richest signal, hosts the largest and cleanest
significant losses ($-9.5\%$ to $-7.3\%$ NDCG@10/Recall@10/MRR@10
across Steam-3k Mul and Steam-15k Mul/Val, all robust).

The architecture was designed with an explicit fallback
($\lambda{=}0$ recovers SASRec), so the failure has a sharp
mechanism question attached, which we answer through an analysis of
the learned per-layer $\lambda$. The learned $\lambda$ does adapt: it
drops to $0.07$ on ML-1M Val (and the model recovers SASRec), so the
optimisation path is open. But on Steam it stays at $0.55$--$0.99$
across all three variants and all three subset sizes, even as test
metrics degrade. The optimiser can disable intensity and chooses not to. A
single per-layer scalar is not a sufficient fallback for
intensity injection: it cannot down-weight specific intensity values,
and it cannot let some attention heads use intensity while others
ignore it. It also confounds the signal itself with the way we encoded
it: a flat $\lambda$ could mean intensity carries no information, or
that our normalisation choice flattened the information that was there.

\paragraph{Contributions.}
\begin{itemize}
\item Three IA-SASRec variants (Add / Mul / Val) realised as a
unified framework with one learnable scalar $\lambda$ per layer, so
that $\lambda{=}0$ collapses each variant exactly to SASRec.
\item A 7-dataset $\times$ 5-seed evaluation under a dual-criterion
significance protocol (paired-$t$ at the seed level and a per-user
paired bootstrap with $B{=}2{,}000$, guarded by a per-seed CI
sign-majority on the bootstrap) explicitly designed against
the failure modes flagged by prior reproducibility critiques~\citep{ferrari2019,sun2020daisy}.
\item A negative result with mechanism evidence: learned $\lambda$
stays large precisely where the variant hurts the most, and the
``off-switch at $\lambda{=}0$'' fallback is not realised in
practice.
\end{itemize}
